\chapter{Conclusion and Discussion}
\section{Summary}
This project investigated the detection of money mule accounts using behavioural
transaction features together with both unsupervised and supervised machine learning
approaches.

A set of account-level features was designed to capture characteristics commonly
associated with money mule behaviour, including transaction frequency, pass-through
activity, draining behaviour, and many-out transfer patterns. These features were first
used in an unsupervised setting through DBSCAN and autoencoder models to identify
accounts whose behaviour deviated from the majority of retail customers. Subsequently,
confirmed mule accounts provided by the bank were incorporated to construct supervised
baselines using logistic regression and XGBoost.

The experimental results suggest that behavioural features provide meaningful
information for distinguishing suspicious accounts. DBSCAN successfully separated a
small set of behavioural outliers, while the autoencoder trained on the presumed normal
population achieved substantially better reconstruction quality than the model trained on
the entire dataset. This supports the assumption that learning only normal behaviour is
more suitable for reconstruction-based anomaly detection. Among the supervised models,
XGBoost produced the strongest overall detection performance, indicating that nonlinear
relationships between behavioural features are important for identifying money mule
activity. Compared with the unsupervised ones, supervised methods enhanced the explainability and cross-validate the feature importance. All methods shown the drain ratio is a very strong feature.

Overall, the study demonstrates that combining domain-driven behavioural feature
engineering with machine learning provides a practical framework for prioritising
potential money mule accounts for further investigation.
\section{Limitation}
Although the proposed framework produced encouraging results, several limitations
remain.

First, computational constraints limited the scale of several experiments. In particular,
the FIFO-based drain ratio feature requires sequential allocation of outgoing
transactions to incoming transactions, making it computationally expensive for large
transaction datasets. Consequently, this feature was evaluated only on a development
dataset containing approximately 2,000 retail accounts rather than the full customer
population as it failed to be added to 20000 accounts due to capacity limit.

Second, the study relied primarily on account-level behavioural features. While these
features capture temporal transaction patterns, they do not explicitly model relationships
between accounts. Previous research has shown that graph-based features, such as
shared counterparties, transaction chains, and network centrality measures, can provide
additional information for detecting organised money mule networks. These features
were not included because constructing large-scale transaction graphs would require
substantially greater computational resources and engineering effort.

Another important limitation concerns the availability of labelled data.
Only a relatively small number of confirmed money mule accounts were
available, while the remaining accounts were unlabeled rather than
verified non-mule accounts. Nevertheless, for the supervised learning
experiments, all unlabeled accounts were treated as negative examples.

This assumption may not hold in practice, since some genuine money mule
accounts may remain unidentified within the unlabeled population.
Consequently, the supervised models were trained using noisy labels,
which may reduce both model performance and the reliability of the
estimated class probabilities. Positive-unlabeled (PU) learning provides
a potential solution to this problem by explicitly accounting for the
presence of hidden positives within the unlabeled data and therefore
represents an important direction for future work.

In addition, the engineered behavioural features were intentionally designed to remain
interpretable and computationally feasible. As a result, some potentially informative
characteristics, including graph-based features like counterparty behaviour, in and out degree, long-range temporal
dependencies, and static features like customer profile information, were not incorporated. These may provide additional discriminatory power when combined with the existing behavioural
features.
\section{Future Work}
Several directions for future work can be considered. 

First, the proposed behavioural
features could be computed on the complete transaction database using distributed
implementations to evaluate scalability in a production environment. In particular, the FIFO-based drain ratio should be redesigned or optimised before being applied to larger samples, since the sequential allocation of outgoing transactions to incoming funds is computationally expensive.

A second direction is to incorporate richer relational information between accounts. The behavioural features developed in this project primarily describe the transaction patterns of individual accounts, while relationships between connected accounts were not explicitly modelled. Counterparty-based features, such as repeated transfers between the same accounts or measures of counterparty diversity, could provide additional discriminatory power with relatively low computational overhead. More sophisticated graph-based approaches, including network centrality measures, community detection and Graph Neural Networks, could further capture the structure of money mule networks and improve detection performance.

Third, more advanced semi-supervised approaches, such as positive-unlabeled (PU) learning
or weakly-supervised representation learning, may better exploit the large amount of
unlabeled transaction data available in banking applications. 

Another promising direction is to incorporate temporal modelling of transaction behaviour. The behavioural features developed in this study are aggregated over a predefined observation period and therefore do not explicitly capture how account behaviour evolves over time. Sequential models, such as recurrent neural networks or Transformer-based architectures, could model the temporal ordering of transactions and detect behavioural changes as they occur. Such approaches may improve the detection of money mules whose transaction patterns develop progressively rather than appearing anomalous only after aggregation.

Finally, combining
behavioural anomaly detection with supervised classification in a hybrid detection
framework may further improve both detection performance and model interpretability.